\subsection{Особенности сходимости вблизи критической точки}

Наиболее сложные численные режимы
возникают вблизи критической точки
и на примыкающей к ней двухфазной границе.
В двухфазной области (регион~4)
температура однозначно задаёт давление насыщения:
\begin{equation}
p = p_{\mathrm{sat}}(T).
\end{equation}
Следовательно, при задании входной пары $(\rho, T)$
дополнительной неизвестной становится не давление,
а степень сухости $x$.

Связь между заданной плотностью смеси
и степенью сухости удобно записывать
через удельный объём.
Для насыщенной жидкости и насыщенного пара
обозначим свойства через штрихи:
$v'$ и $v''$, либо эквивалентно $\rho'$ и $\rho''$.
Тогда для двухфазной смеси выполняется соотношение
\begin{equation}
\begin{aligned}
v &= (1-x)\,v' + x\,v'', \\
v &= \frac{1}{\rho_{\mathrm{target}}}.
\end{aligned}
\end{equation}
Эквивалентная запись через плотности имеет вид
\begin{equation}
\frac{1}{\rho_{\mathrm{target}}}
=
\frac{1-x}{\rho'}
+
\frac{x}{\rho''}.
\end{equation}
Отсюда степень сухости определяется явно:
\begin{equation}
\begin{aligned}
x &= \frac{v - v'}{v'' - v'} \\
&= \frac{\frac{1}{\rho_{\mathrm{target}}} - \frac{1}{\rho'}}
{\frac{1}{\rho''} - \frac{1}{\rho'}}.
\end{aligned}
\end{equation}
Именно это соотношение
связывает заданную среднюю плотность смеси
с массовой долей пара.

После вычисления $x$
аддитивные по массе свойства
определяются по обычному правилу смешения:
\begin{equation}
y = (1-x)\,y' + x\,y''.
\end{equation}
Если вычисленное значение $x$
выходит за диапазон $[0; 1]$,
то состояние не принадлежит двухфазной области
и должно быть возвращено на однофазную маршрутизацию.

Вблизи критической точки
разность $v''-v'$ стремится к нулю.
Поэтому формула для $x$
становится плохо обусловленной:
небольшая ошибка в плотности
даёт заметную ошибку в степени сухости.
В реализации это компенсируется
контролем диапазонов,
проверкой близости к критической точке
и отдельной обработкой пограничных случаев.

